%%%PAGE 192 rot=0 legibility=good summary=光电效应(含电路图)与原子物理要点；四星系统；恒星S2绕黑洞(算黑洞质量)；右侧页边波干涉草算
%%%BLOCK id=192-01 ch=17 sec=光电效应 kind=knowledge alt=none cont=none y=0.04-0.34
1.\ 发生光电效应\ $h\nu>W_0$\quad 跃迁\ $E_{nm}=E_n-E_m$\ 正好\quad 电离\ $E>E_m-E_n$ % 原稿跃迁式中E下方的下标难辨(nm/mn)

$E_k=h\nu-W_0$\hfill 使$R\uparrow$\ $U_{\text{遏止}}\uparrow$ % 右侧“使R↑ U遏止↑”为写在该行右上方的批注
%FIG p192_a 0.085 0.138 0.292 0.285
\notefig[0.25]{figures/p192_a.png}

黄光照\textcircled{A}有示数\quad $\nu_{\text{黄光}}>\nu_{\text{极限}}$ % 原稿此处符号写得形似U/ν，按频率ν转录

增加光强\ \textcircled{A}变大\qquad 该装置可测得从阴极板上逸出的光电子个数

若移动$P$电流表示数可能不变（可能已达饱和光电流） % “移动”二字原稿为写在“若”上方的小字

将$P$移至最左侧，阴阳极板间电压为0，阴极逸出的光电子部分能到达阳极，\textcircled{A}不为0。在某个范围内无光电流（达遏止电压时）。

光强大（频率一定时）$\to$光子数目多$\to$发射出电子多$\to$饱和光电流大

光子频率高$\to$光子能量大$\to$光电子最大初动能大。
%%%END
%%%BLOCK id=192-02 ch=17 sec=射线·γ射线与中子 kind=knowledge alt=none cont=none y=0.345-0.375
2.\ $\gamma$射线与$\nuc{1}{0}{n}$不会发生电磁偏转。（因$\gamma$与$\nuc{1}{0}{n}$不带电）
%%%END
%%%BLOCK id=192-03 ch=17 sec=氢原子能级与光谱 kind=knowledge alt=none cont=none y=0.375-0.405
3.\ 高能级$\to$低能级\ 辐射发出光子，低能级$\to$高能级\ 吸收光子。
%%%END
%%%BLOCK id=192-04 ch=17 sec=核能·质量亏损 kind=knowledge alt=none cont=none y=0.405-0.43
4.\ \underline{原子核静质量}小于\underline{全部核子独存时的总静质量}。 % 原稿“小于”被圈起
%%%END
%%%BLOCK id=192-05 ch=17 sec=波粒二象性·不确定性关系 kind=knowledge alt=none cont=none y=0.425-0.455
5.\ $\Delta x\,\Delta p\ge$\unclear{$\dfrac{h}{4\pi}$}\ 不可能同时确定微观粒子的位置和动量。 % 原稿“位置和动量”被圈起；“≥”后的式子被笔画覆盖难辨，疑为h/4π；此行有一道横线贯穿
%%%END
%%%BLOCK id=192-06 ch=17 sec=原子核·组成 kind=knowledge alt=none cont=none y=0.455-0.495
\underline{原子核不存在独立的电子}（仅有质子中子）
%%%END
%%%BLOCK id=192-07 ch=17 sec=波粒二象性·波动性 kind=knowledge alt=none cont=none y=0.455-0.495
光的波动性\fbox{不是}光子相互作用造成的。
%%%END
%%%BLOCK id=192-08 ch=05 sec=多星系统·四星 kind=problem alt=none cont=none y=0.50-0.70
6.
%FIG p192_b 0.05 0.515 0.18 0.592
\notefig[0.25]{figures/p192_b.png}

\begin{problem}
$G$\ 四星系统
\begin{solution}
\[
d=\frac{\sqrt3}{3}L
\]
做圆周运动的半径\ $r=d=\dfrac{\sqrt3}{3}L$

顶点处\ $F_{\text{向}}=2\dfrac{Gm^2}{L^2}\cos30^\circ+\dfrac{Gm^2}{r^2}=\dfrac{(3+\sqrt3)Gm^2}{L^2}$

由牛二\ $F_{\text{向}}=m\omega^2r$\quad $\omega=\sqrt{\dfrac{(3+3\sqrt3)Gm}{L^3}}$\quad $\checkmark$

$T=\dfrac{2\pi}{\omega}=2\pi\sqrt{\dfrac{L^3}{(3+3\sqrt3)Gm}}$\quad $\checkmark$
\end{solution}
\end{problem}
%%%END
%%%BLOCK id=192-09 ch=05 sec=黑洞·恒星S2绕行 kind=problem alt=none cont=none y=0.70-0.99
7.
\begin{problem}
恒星$S_2$绕黑洞一周16年，离黑洞中心的最近距离为120个天文单位，椭圆半长轴为970天文单位，（地球绕太阳圆周运动的轨迹半径为1天文单位）。恒星$S_2$最高运行速率为$7000\unit{km/s}$。 % 原稿“120”上方标有小字d_1；“120”首位的“1”笔画较细
\begin{solution}
\[
v_{\text{近}}r_{\text{近}}=v_{\text{远}}r_{\text{远}}\qquad r_{\text{远}}=1820\text{天文单位}=(2a-d_1) % 原稿“r远=1820天文单位=(2a-d1)”被一个椭圆圈起
\]
算不了恒星$S_2$质量
%FIG p192_e 0.325 0.812 0.612 0.925
\notefig[0.3]{figures/p192_e.png}
\[
\frac{GM_{\text{黑}}m_{\text{恒}}}{r^2}=m_{\text{恒}}\frac{4\pi^2}{T^2}r
\]
\[
\frac{GM_{\text{日}}M_{\text{地}}}{r_{\text{地}}^2}=m_{\text{地}}\frac{4\pi^2}{T^2}r_{\text{地}}
\]
\[
\frac{M_{\text{黑}}}{M_{\text{日}}}=\frac{r^3T_{\text{地}}^2}{r_{\text{地}}^3T^2}=\frac{970^3\cdot1^2}{1^3\times16^2}=3.6\times10^6
\]
\[
a_{\text{近}}=\frac{GM_{\text{黑}}}{r_{\text{近}}^2}\qquad a_{\text{地}}=\frac{GM_{\text{日}}}{r_{\text{地}}^2}\qquad \frac{a_{\text{近}}}{a_{\text{地}}}=248
\]
\end{solution}
\end{problem}
%%%END
%%%BLOCK id=192-10 ch=08 sec=波的干涉·页边草算 kind=derivation alt=none cont=none y=0.32-0.70
% 页面右侧空白处的草算/草图(无题干；疑为两波源干涉：λ=20、T=1.5，“减”疑指减弱)；另在左下角(四星系统下方)有一幅无标注的曲线草图，疑与之有关
%FIG p192_d 0.71 0.325 0.985 0.445
\notefig[0.3]{figures/p192_d.png}
%FIG p192_c 0.02 0.595 0.145 0.697
\notefig[0.25]{figures/p192_c.png}

$\lambda=20$ % 写在上图下方偏右

$60-2x=10$\ \unclear{20m}

$2x-60=0$

\begin{itemize}
\item 竖排一列(左侧标\unclear{$L$})：10，20，40
\item 竖排一列(右侧)：10，20，40，60
\item 大括号内：$x=30$，$x=10$，$x=0$
\item 另有零散：$x=\unclear{10}$，$x=20$，$x=\unclear{25}$ % CHECK: 页边草算，排布及数字难以还原
\end{itemize}
%%%END
%%%PAGEEND 192
